\lysh{15791_SOLUTION}{1}
α) Ο κύκλος $C_1$ έχει κέντρο $A(0,7)$ και ακτίνα $\rho=2$, άρα αντικαθιστώντας στον τύπο $(x-x_0)^2+(y-y_0)^2=\rho^2$, τις αντίστοιχες τιμές έχουμε:
\[
(x-0)^2+(y-7)^2=2^2 \Leftrightarrow x^2+(y-7)^2=4
\]
β)
i. Σύμφωνα με το τύπο $(x-x_0)^2+(y-y_0)^2=\rho^2$, οι κύκλοι με κέντρο $B(x_1,y_1)$ και ακτίνα $2$ έχουν εξίσωση $(x-x_1)^2+(y-y_1)^2=4$.

ii. Η απόσταση δύο σημείων $A(x_A,y_A)$ και $B(x_B,y_B)$ δίνεται από τον τύπο $AB=\sqrt{(x_A-x_B)^2+(y_A-y_B)^2}$, οπότε για τα σημεία $A(0,7)$ και $B(x_1,y_1)$ έχουμε
\[
AB=\sqrt{(0-x_1)^2+(7-y_1)^2} \Leftrightarrow AB=\sqrt{x_1^2+(7-y_1)^2} \text{.}
\]
γ) Δύο κύκλοι εφάπτονται εξωτερικά αν και μόνο αν η διάκεντρος είναι ίση με το άθροισμα των ακτίνων τους. Οπότε έχουμε:
\[
AB=2+2 \Leftrightarrow \sqrt{(x_1-0)^2+(y_1-7)^2}=4 \Leftrightarrow x_1^2+(y_1-7)^2=16 \quad (1).
\]
Ένας κύκλος εφάπτεται σε ευθεία αν και μόνο αν το κέντρο του κύκλου απέχει από την ευθεία απόσταση ίση με την ακτίνα του. Οπότε έχουμε:
\[
d(B,\varepsilon)=\frac{|0x_1+1y_1-5|}{\sqrt{0^2+1^2}}=2 \Leftrightarrow |y_1-5|=2 \quad (2).
\]
Για να βρούμε τους κύκλους που εφάπτονται στον κύκλο $C_1$ και στην ευθεία $(\varepsilon)$ επιλύουμε το σύστημα των εξισώσεων (1) και (2):
\[
\left\{
\begin{array}{l}
x_1^2+(y_1-7)^2=16 \\
|y_1-5|=2
\end{array}
\right.
\]
\[
\left\{
\begin{array}{l}
x_1^2+(y_1-7)^2=16 \\
x_1-5=2
\end{array}
\right. \Leftrightarrow \left\{
\begin{array}{l}
7^2+(y_1-7)^2=16 \\
x_1=7
\end{array}
\right. \Leftrightarrow \left\{
\begin{array}{l}
(y_1-7)^2=-33 \\
x_1=7
\end{array}
\right. \quad \text{Αδύνατο.}
\]
\[
\text{ή}
\]
\[
\left\{
\begin{array}{l}
x_1^2+(y_1-7)^2=16 \\
x_1-5=-2
\end{array}
\right. \Leftrightarrow \left\{
\begin{array}{l}
3^2+(y_1-7)^2=16 \\
x_1=3
\end{array}
\right. \Leftrightarrow \left\{
\begin{array}{l}
(y_1-7)^2=7 \\
x_1=3
\end{array}
\right. \Leftrightarrow \left\{
\begin{array}{l}
y_1-7=\pm\sqrt{7} \text{ ή } y_1=7+\sqrt{7} \\
x_1=3
\end{array}
\right.
\]
Τελικά είναι δύο οι δύο κύκλοι που εφάπτονται εξωτερικά στον κύκλο $C_1$ και στην ευθεία $(\varepsilon)$ έχουν κέντρα τα σημεία $B(3,7-\sqrt{7})$ και $B'(3,7+\sqrt{7})$ και ακτίνα $\rho=2$.
% TODO: TikZ σχήμα (μην το κατασκευάσεις)